\documentclass[12pt]{article}

\usepackage[T1]{fontenc}
\usepackage[utf8]{inputenc}
\usepackage[romanian]{babel}

\usepackage{amsmath, amssymb, amsthm}
\usepackage{mathrsfs}
\usepackage{geometry}
\geometry{a4paper, margin=2.5cm}

\title{Variabile Continue cu Densitate}
\author{}
\date{}

\theoremstyle{plain}
\newtheorem{definition}{Definiție}
\newtheorem{proposition}{Propoziție}
\newtheorem{example}{Exemplu}

\begin{document}

\maketitle

\section{Variabile aleatoare continue}

\begin{definition}
O variabilă aleatoare $X$ se numește \textbf{continuă} dacă există o funcție
$f : \mathbb{R} \to [0,\infty)$, numită \textbf{densitate}, astfel încât:


\[
\mathbb{P}(X \in B) = \int_B f(x)\, dx
\]


pentru orice mulțime boreliană $B \subseteq \mathbb{R}$.
Funcția $f$ satisface:


\[
f(x) \ge 0, \qquad \int_{-\infty}^{\infty} f(x)\, dx = 1.
\]


\end{definition}

\section{Exemple fundamentale de distribuții continue}

\subsection{Distribuția uniformă pe un interval}

\begin{definition}
Fie $a < b$. O variabilă aleatoare $X$ are distribuție \textbf{uniformă} pe intervalul $[a,b]$ dacă:


\[
f(x) =
\begin{cases}
\dfrac{1}{b-a}, & x \in [a,b], \\[12pt]
0, & \text{altfel}.
\end{cases}
\]


Notăm $X \sim \mathrm{Unif}(a,b)$.
\end{definition}

\subsection{Distribuția exponențială}

\begin{definition}
O variabilă aleatoare $X$ se numește \textbf{exponențială} cu parametru $\lambda > 0$ dacă are densitatea:


\[
f(x) =
\begin{cases}
\lambda e^{-\lambda x}, & x \ge 0,\\
0, & x < 0.
\end{cases}
\]


Notăm $X \sim \mathrm{Exp}(\lambda)$.
\end{definition}

\subsection{Distribuția Gamma}

\begin{definition}
Fie $\alpha > 0$ (parametrul de formă) și $\lambda > 0$ (parametrul de rată).
O variabilă aleatoare $X$ are distribuție \textbf{Gamma} dacă:


\[
f(x) =
\begin{cases}
\dfrac{\lambda^\alpha}{\Gamma(\alpha)} x^{\alpha - 1} e^{-\lambda x}, & x \ge 0,\\[6pt]
0, & x < 0.
\end{cases}
\]


Notăm $X \sim \mathrm{Gamma}(\alpha,\lambda)$.
\end{definition}

\subsection{Distribuția normală (gaussiană)}

\begin{definition}
Fie $\mu \in \mathbb{R}$ și $\sigma > 0$.
O variabilă aleatoare $X$ se numește \textbf{normală} cu medie $\mu$ și deviație standard $\sigma$ dacă:


\[
f(x) = \frac{1}{\sigma \sqrt{2\pi}}
\exp\!\left( -\frac{(x - \mu)^2}{2\sigma^2} \right).
\]


Notăm $X \sim \mathcal{N}(\mu,\sigma^2)$.
\end{definition}

\end{document}
